% Rapport technique - Mangaka : transfert de style photo -> anime
% Compilation : pdflatex rapport_technique.tex (deux fois pour la table des matières)
% Convention : tout ce qui est marqué \AMESURER ou \AFAIRE est à compléter.
% Aucun chiffre ne doit être écrit sans avoir été mesuré (MLflow, models/metrics.json).

\documentclass[11pt,a4paper]{article}
\usepackage[utf8]{inputenc}
\usepackage[T1]{fontenc}
\usepackage[french]{babel}
\usepackage[margin=2.5cm]{geometry}
\usepackage{graphicx}
\usepackage{booktabs}
\usepackage{amsmath}
\usepackage{xcolor}
\usepackage{hyperref}
\usepackage{listings}

% Emplacements à remplir : bien visibles en rouge dans le PDF
\newcommand{\AMESURER}{\textcolor{red}{\textbf{[à mesurer]}}}
\newcommand{\AFAIRE}[1]{\textcolor{red}{\textbf{[À FAIRE : #1]}}}
\newcommand{\FIGURE}[1]{\fbox{\parbox{0.9\linewidth}{\centering\vspace{2cm}
  \textcolor{red}{\textbf{Figure à insérer : #1}}\vspace{2cm}}}}

\title{Mangaka : transfert de style photo vers anime\\
  \large Projet de module Deep Learning Avancé -- Axe 2 (IA générative)}
\author{\AFAIRE{Nom 1} \and \AFAIRE{Nom 2}}
\date{\AFAIRE{date de la soutenance}}

\begin{document}
\maketitle
\tableofcontents
\newpage

% ---------------------------------------------------------------------------
\section{Problématique}

Transformer une photo réelle en image au style d'un réalisateur d'animation (Hayao Miyazaki,
Makoto Shinkai) en conservant son contenu. Il n'existe pas de paires photo/dessin : le
problème est une traduction d'image à image \emph{non appariée}.

Deux exigences contradictoires, mesurées séparément :
\begin{itemize}
  \item proximité au style visé : \textbf{FID} entre images générées et vraies images d'anime ;
  \item préservation du contenu : \textbf{LPIPS} entre la photo et sa version stylisée.
\end{itemize}

Cas d'usage : application web où l'utilisateur envoie une photo (fichier, galerie, webcam),
choisit un style et reçoit le résultat avec la latence d'inférence.

% ---------------------------------------------------------------------------
\section{Données et prétraitement}

\subsection{Jeu de données}
Dataset public d'AnimeGAN (Chen et al., 2019) : photos réelles d'entraînement
(\texttt{train\_photo}), photos de test (\texttt{test\_photo}) et images extraites de films
(\texttt{Hayao/style}, \texttt{Shinkai/style}).

\begin{table}[h]
\centering
\begin{tabular}{lr}
\toprule
Sous-ensemble & Nombre d'images \\
\midrule
Photos d'entraînement & 6\,656 \\
Photos de test & 791 \\
Style Hayao & 1\,792 \\
Style Shinkai & 1\,650 \\
\bottomrule
\end{tabular}
\caption{Taille du jeu de données (versionné par DVC, \texttt{data/raw.dvc}).}
\end{table}

\subsection{Prétraitement}
\begin{itemize}
  \item Redimensionnement (256$\times$256 pour AnimeGAN, 128$\times$128 pour CycleGAN) et
        normalisation dans $[-1,1]$, même convention que l'API (\texttt{src/data/dataset.py}).
  \item Lissage des contours des images d'anime (Canny + dilatation + flou gaussien) : le
        discriminateur apprend qu'un trait flou est « faux », ce qui force des contours nets.
  \item Augmentation : miroir horizontal des photos, à l'entraînement uniquement.
\end{itemize}

\subsection{Absence de fuite de données}
Les photos de test ne servent qu'à \texttt{evaluate.py}. \AFAIRE{résultat du contrôle de
doublons entre \texttt{train\_photo} et \texttt{test\_photo} (empreintes des fichiers)}.
Aucune statistique n'est calculée sur l'ensemble de test avant l'évaluation.

% ---------------------------------------------------------------------------
\section{Architectures}

\subsection{Baseline : CycleGAN (128$\times$128)}
Deux générateurs ResNet (6 blocs résiduels, normalisation d'instance) et deux discriminateurs
PatchGAN. Pertes : adverse LSGAN, cycle ($\lambda=10$), identité ($\lambda=5$), historique de
50 images générées pour le discriminateur (\texttt{src/models/train\_cyclegan.py}).
\[
\mathcal{L}_G = \mathcal{L}_{adv} + \lambda_{cyc}\,\lVert G_{ap}(G_{pa}(x)) - x\rVert_1
  + \lambda_{id}\,\lVert G_{pa}(y) - y\rVert_1 + \dots
\]

\subsection{Modèle avancé : AnimeGAN (256$\times$256)}
Générateur léger de type AnimeGANv2 (blocs résiduels inversés, convolutions en profondeur),
un générateur par style. Discriminateur PatchGAN avec \textbf{normalisation spectrale}.
Pertes : adverse, contenu (VGG19), style (matrices de Gram sur images en gris), couleur (YUV).
Phase d'initialisation (reconstruction seule) avant le jeu adverse
(\texttt{src/models/train.py}). Le style Shinkai est affiné depuis les poids Hayao
(\texttt{--resume}).

\begin{table}[h]
\centering
\begin{tabular}{lcc}
\toprule
 & CycleGAN & AnimeGAN \\
\midrule
Paramètres du générateur & 7\,837\,699 & 973\,699 \\
Taille du fichier ONNX & 30,0 Mo & 3,8 Mo \\
\bottomrule
\end{tabular}
\caption{Taille des modèles.}
\end{table}

% ---------------------------------------------------------------------------
\section{Protocole expérimental}
\begin{itemize}
  \item Matériel : GPU Kaggle 2 $	imes$ NVIDIA Tesla T4, durée d'entraînement \AMESURER.
  \item Hyperparamètres : voir les paramètres journalisés dans MLflow (époques, lr, poids des
        pertes) \AFAIRE{tableau des valeurs retenues}.
  \item Évaluation sur \texttt{test\_photo} uniquement : FID (Inception v3, 2048
        caractéristiques) contre les images du style, LPIPS (AlexNet) contre la photo d'entrée.
  \item Export ONNX (opset 17, hauteur et largeur dynamiques) et vérification de l'écart
        PyTorch / ONNX Runtime sur la même image.
\end{itemize}

% ---------------------------------------------------------------------------
\section{Résultats : baseline contre modèle avancé}

\begin{table}[h]
\centering
\begin{tabular}{lcccc}
\toprule
Modèle & Résolution & FID $\downarrow$ & LPIPS $\downarrow$ & Latence ONNX CPU (ms) \\
\midrule
CycleGAN (Hayao) & 128 & 144,1 & 0,128 & 168,7 \\
AnimeGAN (Hayao) & 256 & \AMESURER & \AMESURER & \AMESURER \\
AnimeGAN (Shinkai) & 256 & \AMESURER & \AMESURER & \AMESURER \\
\bottomrule
\end{tabular}
\caption{Comparaison sur les photos de test (source : \texttt{models/metrics.json}).}
\end{table}

\begin{table}[h]
\centering
\begin{tabular}{lcc}
\toprule
 & PyTorch CPU & ONNX Runtime CPU \\
\midrule
Latence moyenne (ms) & \AMESURER & \AMESURER \\
Écart maximal des sorties & \multicolumn{2}{c}{\AMESURER} \\
\bottomrule
\end{tabular}
\caption{Gain de l'export ONNX.}
\end{table}

\FIGURE{courbes de pertes générateur / discriminateur (captures MLflow)}

\FIGURE{grille photo / CycleGAN / AnimeGAN Hayao / AnimeGAN Shinkai sur 4 photos de test}

% ---------------------------------------------------------------------------
\section{Analyse critique}
\AFAIRE{à rédiger à partir des résultats mesurés}
\begin{itemize}
  \item Compromis FID / LPIPS : un style plus marqué éloigne du contenu.
  \item Effet de la résolution (128 contre 256) sur la comparaison.
  \item Artefacts observés (couleurs, contours, textures) et leur cause probable.
  \item Effet de la normalisation spectrale sur la stabilité de l'entraînement.
\end{itemize}

% ---------------------------------------------------------------------------
\section{MLOps}
\begin{itemize}
  \item \textbf{DVC} : \texttt{data/raw} versionné (\texttt{data/raw.dvc}), remote
        \AFAIRE{type de remote} ; pipeline \texttt{dvc.yaml} prepare $\rightarrow$ train
        $\rightarrow$ evaluate.
  \item \textbf{MLflow} : paramètres, pertes par époque, grilles d'images par époque,
        checkpoints ; Model Registry avec statut Staging puis Production
        (\texttt{src/models/registry.py}).
  \item \textbf{API} : FastAPI asynchrone, schémas Pydantic, codes 400 / 422 / 503, ONNX
        Runtime, modèle factice tant qu'aucun \texttt{.onnx} n'est présent.
  \item \textbf{Frontend} : React + Tailwind, glisser-déposer, galerie, webcam (capture et
        continu), curseur avant/après, latences client et serveur.
  \item \textbf{Docker} : Dockerfiles multi-stage, \texttt{docker compose up --build}
        (frontend, backend, MLflow).
  \item \textbf{CI} : GitHub Actions (flake8, black, pytest, eslint, build Vite, build des
        images Docker et test de \texttt{/health}).
\end{itemize}

\FIGURE{capture de l'interface MLflow (runs et Model Registry)}

% ---------------------------------------------------------------------------
\section{Limites et perspectives}
\begin{itemize}
  \item Pas d'espace latent pour un générateur feed-forward : \AFAIRE{décision prise sur les
        curseurs d'interpolation demandés par le sujet}.
  \item Budget GPU limité (Kaggle) : nombre d'époques réduit, poids des pertes peu réglés.
  \item FID calculé sur un petit nombre d'images de test : estimation bruitée.
  \item Inférence sur CPU dans Docker ; quantification INT8 non explorée.
\end{itemize}

\end{document}
